% ── §3 Karnaugh Maps and Quine–McCluskey Minimisation ──
\section{Karnaugh Maps and Quine--McCluskey Minimisation}
\label{sec:theory-kmap-qm}

\subsection{Karnaugh Maps for Residue-Pair Statistics}
\label{sec:kmap-construction}

For a pair of alignment columns $(i, j)$ with $i < j$, we define a ternary
Karnaugh map $K_{ij}$ as a $32 \times 32$ grid whose rows and columns are
indexed by the five-bit Gray codes of the consensus residues at positions $i$
and $j$, respectively.  Each cell $(a, b) \in \{0,\ldots,31\}^2$ of the map
takes one of three values:
\begin{equation}
	K_{ij}(a, b) =
	\begin{cases}
		+1 & \text{if the residue pair } (a, b) \text{ is an observed mutation}, \\
		-1 & \text{if } (a, b) \text{ is the reference (majority) pair},         \\
		0  & \text{if } (a, b) \text{ has never been observed}.
	\end{cases}
	\label{eq:kmap-def}
\end{equation}
Rows and columns with indices $20$--$31$ correspond to unused Gray codewords
and are assigned don't-care status ($-1$), which permits the
Quine--McCluskey algorithm to form larger implicants by merging on-set cells
through these padding regions.  This padding step constitutes the second
critical correction applied to our pipeline; without it, a $20 \times 20$
map fed to a minimiser expecting $4^k = 2^{10} = 1024$ cells would silently
wrap indices $256$--$399$ onto indices $0$--$143$, corrupting minterm labels
and producing phantom rules.

\subsection{Ternary Truth Tables}
\label{sec:truth-tables}

For computational purposes, each Karnaugh map is flattened into a one-dimensional
truth table $\tau : \{0, \ldots, 1023\} \to \{-1, 0, 1\}$ where the cell index
$c$ encodes both residue identities in MSB-first bit order:
\begin{equation}
	c = 32 \cdot r + s, \qquad r = g(\mathrm{code}(a_i)), \quad s = g(\mathrm{code}(a_j)),
	\label{eq:cell-index}
\end{equation}
and the bit at position $b$ ($b = 0, \ldots, 9$) of cell $c$ is extracted as
$(c \gg (9 - b)) \mathbin{\&} 1$.  An \emph{implicant} over $w$ variables is a
pair $(\mathrm{values}, \mathrm{mask}) \in \{0,1\}^w \times \{0,1\}^w$ where a
mask bit equal to~1 marks the corresponding variable position as don't-care.
An implicant \emph{matches} a cell $c$ when every constrained (mask = 0)
position agrees with the corresponding bit of $c$:
\begin{equation}
	\mathrm{matches}(v, m, c) \iff
	\bigl((v \oplus c) \mathbin{\&} (\sim m)\bigr) = 0,
	\label{eq:implicant-match}
\end{equation}
where $\sim m$ denotes the bitwise complement of $m$ within the $w$-bit field.

\subsection{Quine--McCluskey Algorithm}
\label{sec:qm-algorithm}

The Quine--McCluskey algorithm~\cite{quine1952,mccluskey1956} finds all prime
implicants of a Boolean function by exhaustive pairwise merging of minterms
that differ in exactly one bit position.  Our implementation extends the
classical formulation to handle ternary truth tables by including don't-care
cells in the initial minterm set alongside on-set cells, thereby allowing
prime implicants to cover don't-care regions without requiring them to do so.
The merge phase groups minterms by popcount of their fixed bits and restricts
comparisons to adjacent popcount groups, reducing the worst-case complexity
from $O(m^2)$ to $O(m \cdot |\text{groups}|)$ for $m$ minterms.

After all merges are exhausted, the surviving implicants that cannot be merged
further are the \emph{prime implicants}.  Among these, a prime implicant is
\emph{essential} if it covers at least one on-set cell that no other prime
implicant covers.  A minimum cover is then constructed greedily from the
essential prime implicants plus additional non-essential primes selected to
cover any remaining uncovered on-set cells.  Two properties of this procedure
are machine-verified in Lean~4:

\begin{theorem}[Cover Completeness]
	\label{thm:cover-complete}
	Every on-set cell of a labeled Karnaugh map is matched by at least one prime
	implicant in the Quine--McCluskey output.
\end{theorem}

\begin{theorem}[Cover Soundness]
	\label{thm:cover-sound}
	No prime implicant matches any off-set cell of a labeled Karnaugh map.
\end{theorem}

Together, these two properties guarantee that the SOP expression produced by
the minimiser neither omits nor overapproximates the coevolutionary signal,
provided the input truth table correctly distinguishes observed mutations
from unobserved combinations and from reference-pair don't-cares.

\subsection{Flipped Maps and Negative Selection}
\label{sec:flipped-maps}

In addition to the positive (mutation-detecting) map described above, we
construct a \emph{flipped} map for each candidate pair by assigning on-set
status to cells that have never been observed across all sequences in the
alignment, off-set status to cells that have been observed, and don't-care
status to reference pairs.  Prime implicants of the flipped map encode
\emph{forbidden} residue-pair combinations---pairs whose joint absence across
all sequenced homologues suggests destabilising interactions or steric
incompatibility.  Together, the positive and flipped maps partition the
candidate space into three disjoint regions: coevolutionary constraints,
negative-selection boundaries, and uninformative gaps.
